% GENERATED FILE. Edit canonical lessons, then run npm run generate:books.
% canonical-source-hash: fnv1a64:1b3cdcdf01c23bc3
% canonical-lessons: JA-C210-ki-tree, JA-C210-ueru, JA-C210-shokubutsu, JA-C210-sodatsu, JA-C210-kareru

\chapter{Trees and Growing Things}
\label{ch:ja-trees-and-growing-things}
\hlchaptermodality{210}

\begin{chapteropening}
\textbf{By the end of this chapter:} \emph{I can say what I will plant and how it grows.}

Complete the last lesson of chapter 210: I can say what I will plant and how it grows.
\end{chapteropening}

\section[\texorpdfstring{\ja{き} (ki)}{き (ki)}]{\ja{き} (ki) — a tree; wood}

\label{lesson:JA-C210-ki-tree}



\begin{quote}
Before the new one: say the Japanese for a cherry tree, then the Japanese for to scatter.
\end{quote}

\subsection*{You'll want to know: \ja{き}}
\textbf{\ja{き}} — \emph{ki} — \textquotedblleft{}a tree; wood\textquotedblright{}.

\begin{quote}
\textbf{\ja{やまには} \ja{きが} \ja{おおいです。}} — \emph{yama ni wa ki ga ōi desu} — \textquotedblleft{}There are many trees on the mountain.\textquotedblright{}
\end{quote}

\begin{cousinweb}
\textbf{\ja{き}} (\emph{ki}) is a native Japanese word (wago), inherited rather than borrowed from Chinese; its kanji was attached later as a kun'yomi, and in compounds it can take the shape \emph{ko} (JMdict, KANJIDIC2).
\end{cousinweb}

\begin{grammarlens}[title={What you've built}]
One more word for everyday life at A2.
\end{grammarlens}

\subsection*{Your turn}
\begin{itemize}
\raggedright
  \item \emph{Say it:} \emph{ki}
  \item \emph{Say it:} \emph{ki}, once more
  \item \emph{Say it:} \emph{yama ni wa ki ga ōi desu} — There are many trees on the mountain.
  \item \emph{Recall:} say the Japanese for a cherry tree, then the Japanese for to scatter, then say \emph{ki} again
\end{itemize}

\begin{tcolorbox}[breakable,colback=teal!4,colframe=teal!35!black,title={Before you move on}]
What does \ja{き} mean? (\textquotedblleft{}A tree; wood\textquotedblright{}.) Say it once more.
\end{tcolorbox}
\section[\texorpdfstring{\ja{うえる} (ueru)}{うえる (ueru)}]{\ja{うえる} (ueru) — to plant}

\label{lesson:JA-C210-ueru}



\begin{quote}
Before the new one: say the Japanese for to scatter, then the Japanese for a tree.
\end{quote}

\subsection*{You'll want to know: \ja{うえる}}
\textbf{\ja{うえる}} — \emph{ueru} — \textquotedblleft{}to plant\textquotedblright{}.

\begin{quote}
\textbf{\ja{らいしゅう} \ja{にわに} \ja{きを} \ja{うえます。}} — \emph{raishū niwa ni ki o uemasu} — \textquotedblleft{}Next week I will plant a tree in the garden.\textquotedblright{}
\end{quote}

Polite \textbf{\ja{うえます}}, \emph{uemasu}; an ichidan verb.

\begin{cousinweb}
\textbf{\ja{うえる}} (\emph{ueru}) is a native Japanese word (wago), inherited rather than borrowed from Chinese; its kanji was attached later as a kun'yomi, and its older origin is not settled (JMdict, KANJIDIC2).
\end{cousinweb}

\begin{grammarlens}[title={What you've built}]
One more word for everyday life at A2.
\end{grammarlens}

\subsection*{Your turn}
\begin{itemize}
\raggedright
  \item \emph{Say it:} \emph{ueru}
  \item \emph{Say it:} \emph{ueru}, once more
  \item \emph{Say it:} \emph{raishū niwa ni ki o uemasu} — Next week I will plant a tree in the garden.
  \item \emph{Recall:} say the Japanese for to scatter, then the Japanese for a tree, then say \emph{ueru} again
\end{itemize}

\begin{tcolorbox}[breakable,colback=teal!4,colframe=teal!35!black,title={Before you move on}]
What does \ja{うえる} mean? (\textquotedblleft{}To plant\textquotedblright{}.) Say it once more.
\end{tcolorbox}
\section[\texorpdfstring{\ja{しょくぶつ} (shokubutsu)}{しょくぶつ (shokubutsu)}]{\ja{しょくぶつ} (shokubutsu) — a plant}

\label{lesson:JA-C210-shokubutsu}



\begin{quote}
Before the new one: say the Japanese for a tree, then the Japanese for to plant.
\end{quote}

\subsection*{You'll want to know: \ja{しょくぶつ}}
\textbf{\ja{しょくぶつ}} — \emph{shokubutsu} — \textquotedblleft{}a plant\textquotedblright{}.

\begin{quote}
\textbf{\ja{この} \ja{しょくぶつは} \ja{みずが} \ja{ひつようです。}} — \emph{kono shokubutsu wa mizu ga hitsuyō desu} — \textquotedblleft{}This plant needs water.\textquotedblright{}
\end{quote}

\begin{cousinweb}
\textbf{\ja{しょくぶつ}} (\emph{shokubutsu}) is Sino-Japanese (kango): two Chinese characters read with their on'yomi, readings borrowed from Chinese. They are \emph{shoku}, \textquotedblleft{}to plant\textquotedblright{} (Middle Chinese \emph{dzyik}) and \emph{butsu}, \textquotedblleft{}thing\textquotedblright{} (\emph{mjut}) (KANJIDIC2; Middle Chinese in Baxter's transcription).
\end{cousinweb}

\begin{grammarlens}[title={What you've built}]
One more word for everyday life at A2.
\end{grammarlens}

\subsection*{Your turn}
\begin{itemize}
\raggedright
  \item \emph{Say it:} \emph{shokubutsu}
  \item \emph{Say it:} \emph{shokubutsu}, once more
  \item \emph{Say it:} \emph{kono shokubutsu wa mizu ga hitsuyō desu} — This plant needs water.
  \item \emph{Recall:} say the Japanese for a tree, then the Japanese for to plant, then say \emph{shokubutsu} again
\end{itemize}

\begin{tcolorbox}[breakable,colback=teal!4,colframe=teal!35!black,title={Before you move on}]
What does \ja{しょくぶつ} mean? (\textquotedblleft{}A plant\textquotedblright{}.) Say it once more.
\end{tcolorbox}
\section[\texorpdfstring{\ja{そだつ} (sodatsu)}{そだつ (sodatsu)}]{\ja{そだつ} (sodatsu) — to grow; to grow up}

\label{lesson:JA-C210-sodatsu}



\begin{quote}
Before the new one: say the Japanese for to plant, then the Japanese for a plant.
\end{quote}

\subsection*{You'll want to know: \ja{そだつ}}
\textbf{\ja{そだつ}} — \emph{sodatsu} — \textquotedblleft{}to grow; to grow up\textquotedblright{}.

\begin{quote}
\textbf{\ja{この} \ja{きは} \ja{よく} \ja{そだちます。}} — \emph{kono ki wa yoku sodachimasu} — \textquotedblleft{}This tree grows well.\textquotedblright{}
\end{quote}

Polite \textbf{\ja{そだちます}}, \emph{sodachimasu}; a godan verb.

\begin{cousinweb}
\textbf{\ja{そだつ}} (\emph{sodatsu}) is a native Japanese word (wago), inherited rather than borrowed from Chinese; its kanji was attached later as a kun'yomi, and its older origin is not settled (JMdict, KANJIDIC2).
\end{cousinweb}

\begin{grammarlens}[title={What you've built}]
One more word for everyday life at A2.
\end{grammarlens}

\subsection*{Your turn}
\begin{itemize}
\raggedright
  \item \emph{Say it:} \emph{sodatsu}
  \item \emph{Say it:} \emph{sodatsu}, once more
  \item \emph{Say it:} \emph{kono ki wa yoku sodachimasu} — This tree grows well.
  \item \emph{Recall:} say the Japanese for to plant, then the Japanese for a plant, then say \emph{sodatsu} again
\end{itemize}

\begin{tcolorbox}[breakable,colback=teal!4,colframe=teal!35!black,title={Before you move on}]
What does \ja{そだつ} mean? (\textquotedblleft{}To grow; to grow up\textquotedblright{}.) Say it once more.
\end{tcolorbox}
\section[\texorpdfstring{\ja{かれる} (kareru)}{かれる (kareru)}]{\ja{かれる} (kareru) — to wither; to die (of a plant)}

\label{lesson:JA-C210-kareru}



\begin{quote}
Before the new one: say the Japanese for a plant, then the Japanese for to grow.
\end{quote}

\subsection*{You'll want to know: \ja{かれる}}
\textbf{\ja{かれる}} — \emph{kareru} — \textquotedblleft{}to wither; to die (of a plant)\textquotedblright{}.

\begin{quote}
\textbf{\ja{ふゆに} \ja{くさが} \ja{かれます。}} — \emph{fuyu ni kusa ga karemasu} — \textquotedblleft{}The grass withers in winter.\textquotedblright{}
\end{quote}

Polite \textbf{\ja{かれます}}, \emph{karemasu}; an ichidan verb.

\begin{cousinweb}
\textbf{\ja{かれる}} (\emph{kareru}) is a native Japanese word (wago), inherited rather than borrowed from Chinese; its kanji was attached later as a kun'yomi, and its older origin is not settled (JMdict, KANJIDIC2).
\end{cousinweb}

\begin{grammarlens}[title={What you've built}]
One more word for everyday life at A2.
\end{grammarlens}

\subsection*{Your turn}
\begin{itemize}
\raggedright
  \item \emph{Say it:} \emph{kareru}
  \item \emph{Say it:} \emph{kareru}, once more
  \item \emph{Say it:} \emph{fuyu ni kusa ga karemasu} — The grass withers in winter.
  \item \emph{Recall:} say the Japanese for a plant, then the Japanese for to grow, then say \emph{kareru} again
\end{itemize}

\begin{tcolorbox}[breakable,colback=teal!4,colframe=teal!35!black,title={Before you move on}]
What does \ja{かれる} mean? (\textquotedblleft{}To wither; to die (of a plant)\textquotedblright{}.) Say it once more.
\end{tcolorbox}
